¥chapter{UNIX コマンド入門}

¥section*{UNIX のコマンドを使ってみよう}

macOS や WSL 上の Ubuntu などの UNIX 系システムでは、キーボードからコマンドを打ち込んで、プログラムを実行することが多い。
コマンドに慣れるために、この章では UNIX の簡単なコマンドをいくつか紹介しよう。

コマンドは、英語を省略したものが多い。
英語と言っても、日本語の会話にもよく出てくるような、簡単なものがほとんどである。
UNIX 系システムでプログラミングをするくらいであれば、
極端な話 10 個くらいのコマンドを知っていればなんとかなる。

とは言え、コマンドは一度覚えてしまえば、macOS や Ubuntu など多くの UNIX 系システムで、
ほとんど同じように使うことができる。
また、たくさんのコマンドやテクニックを知っていると、簡単な
手順で実にいろいろなことができるようになる。
この本に出てくるコマンドは、全部おぼえてしまって損はないだろう。

¥section{システム情報の表示: ¥texttt{uname} コマンド}

¥subsection{¥texttt{uname} コマンドを使ってみよう}

UNIX 系システムにログインすると、macOS では「¥texttt{username@MacBook-Air ¥~}」、
Ubuntu では「¥texttt{username@hostname:¥~¥$}」のような文字が表示されている。
これを「プロンプト -- prompt(入力促進記号)」と言う。
プロンプトが表示されているのは、これに続けてコマンドが入力できるという
お知らせである。

なお、この本では、プロンプトはmacOSでは「¥texttt{¥%}」、Ubuntuでは「¥texttt{¥$}」と書いておくので注意して欲しい。
プロンプトをどんなものにするかは、自分で自由に設定することができる。

¥bigskip

では、プロンプトに対して、
次のように「¥texttt{uname -a}¥retrn」と打ち込んでみよう。
UNIX 系システムでは、¥underline{大文字と小文字は区別}されるので、
このコマンドは正確に入力すること。

なお、この本では、ユーザが入力する文字は「¥textbf{太字}」で、
そうでない文字は「普通の書体」で表現している。

¥begin{progls}
% ¥slbfl{uname -a}¥retrn
¥end{progls}

すると、次のように、システム情報が表示される。
さっき、UNIX 系システムのコマンドは英語を省略したものが多いと言ったが、
「¥texttt{uname}」は「¥slbfl{u}nix ¥slbfl{name} (UNIXの名前)」を省略したものだ。

macOSの場合:
¥begin{progls}
% ¥slbfl{uname -a}¥retrn
Darwin MacBook-Air.local 24.5.0 Darwin Kernel Version 24.5.0: Thu Apr 24 22:15:47 PDT 2025; root:xnu-10063.121.3~1/RELEASE_ARM64_T8112 arm64
%
¥end{progls}

WSL上のUbuntuの場合:
¥begin{progls}
$ ¥slbfl{uname -a}¥retrn
Linux DESKTOP-ABC123 5.15.133.1-microsoft-standard-WSL2 #1 SMP Thu Apr 18 19:04:30 UTC 2025 x86_64 x86_64 x86_64 GNU/Linux
$
¥end{progls}

¥subsection{コマンドを打ったけどうまくいかない??}

コマンドを入力するときにありがちなトラブルとその解決方法を
紹介しよう。

¥begin{itemize}
¥item ¥textbf{打ち間違えた文字を消すには?}

文字を打ち間違えたときは、 ¥retrn を押す前であれば、
 ¥KEYTOP{Delete} や ¥KEYTOP{Back Space} で消すことができる。

¥item ¥textbf{「Command not found.」って言われたけど?}

「¥texttt{cal}」を、「¥texttt{cak}」のように¥textbf{ミス・タイプ}したりすると、
次のように「Command not found.」
 (コマンドが見つからないよ。) というメッセージがでる。

¥begin{progls}
% ¥slbfl{cak}¥retrn
cak: Command not found. ¥textnormal{(コマンドが見つからないよ。)}
%
¥end{progls}

¥item ¥textbf{コマンドを打ったら暴走した?}

何かのキーを間違って押したり、ミス・タイプすると、運が悪い場合、
プログラムが暴走してプロンプトが出て来なくことがある。
そんなときは、 ¥KEYTOP{Ctrl}+c (¥KEYTOP{Ctrl} を押しながら ¥KEYTOP{c})
と打てば、多くの場合止めることができる。
¥end{itemize}

¥subsection{¥texttt{uname} コマンドの引き数}

もう少し高度な ¥texttt{uname} コマンドの使い方を紹介しよう。

¥texttt{uname} コマンドでは、様々なシステム情報を表示させることができる。
「¥texttt{uname -s}¥retrn」というふうにオプションを指定して実行してみよう。

¥begin{progls}
% ¥slbfl{uname -s}¥retrn
Darwin
%
¥end{progls}

このように、macOSの場合はカーネル名「Darwin」が表示されます。Ubuntuの場合は「Linux」と表示されます。

%¥newpage

次は「-r」オプションを使って、「¥texttt{uname -r}¥retrn」と実行してみよう。

¥begin{progls}
% ¥slbfl{uname -r}¥retrn
24.5.0
%
¥end{progls}

これはカーネルのリリースバージョンを表示しています。実行環境によって表示される数値は異なります。

¥texttt{uname} コマンドには他にも様々なオプションがあります：

macOSの場合:
¥begin{progls}
% ¥slbfl{uname -m}¥retrn
arm64
%
¥end{progls}

WSL上のUbuntuの場合:
¥begin{progls}
$ ¥slbfl{uname -m}¥retrn
x86_64
$
¥end{progls}

これはマシンのハードウェアタイプを表示しています。

¥bigskip

ひとまず、ここまでのところをまとめよう。
¥texttt{uname} コマンドは次のように使う。

¥begin{coml}{5cm}
¥texttt{uname} 「オプション」
¥end{coml}

主なオプションは以下の通り：
\begin{itemize}
\item -a: 全ての情報を表示
\item -s: カーネル名を表示
\item -n: ホスト名を表示
\item -r: カーネルのリリースバージョンを表示
\item -m: マシンのハードウェアタイプを表示
\end{itemize}

オプションを省略すると、カーネル名のみが表示される。

¥bigskip

このようなコマンドのオプション指定は、多くのUNIX系システムコマンドで共通して使われています。
オプションは通常「-」（ハイフン）で始まり、その後に1文字や複数文字が続きます。
これは、コマンドの基本的な動作を変更したり、追加情報を指定するために使われます。
このような指定方法は、
 ¥textbf{「¥ruby{引}{¥mc ひ}き¥ruby{数}{¥mc すう} -- argument」}または¥textbf{「オプション -- option」}と言います。
これはプログラミングの本や、マニュアルなどでもよく使う用語なので、
おぼえておこう。

¥vspace*{6mm}

¥begin{BOXNOTE}
{¥large¥bf ☆練習☆}
¥begin{enumerate}
¥item ¥texttt{uname -a} コマンドを実行し、システムについて表示される全ての情報を確認してみよう。

¥item macOSとUbuntu（またはその他のLinux）で ¥texttt{uname} コマンドの結果を比較して、どのような違いがあるか調べてみよう。
¥end{enumerate}
¥end{BOXNOTE}

¥section{その他の簡単なコマンド}

¥subsection{日時の表示: ¥texttt{date} コマンド}

 ¥texttt{date} は、現在の時刻を表示するコマンドだ。
「¥texttt{date}¥retrn」と打ってみよう。
なお、今後は ¥retrn を一々書かないで省略する。

¥begin{progls}
% ¥slbfl{date}
2025年 4月 25日 金曜日 14時30分15秒 JST
% 
¥end{progls}

また、¥texttt{date} コマンドにはフォーマット指定オプションがあり、日付や時刻の表示形式をカスタマイズできます。
例えば、「¥texttt{date "+%Y年%m月%d日 (%a) %H:%M:%S"}」と実行すると、以下のような形式で表示されます。

¥begin{progls}
% ¥slbfl{date "+%Y年%m月%d日 (%a) %H:%M:%S"}
2025年04月25日 (金) 14:30:15
% 
¥end{progls}

主なフォーマット指定子は以下の通りです：
\begin{itemize}
\item %Y: 西暦（4桁）
\item %m: 月（01～12）
\item %d: 日（01～31）
\item %a: 曜日の略称
\item %H: 時（00～23）
\item %M: 分（00～59）
\item %S: 秒（00～59）
\end{itemize}


¥subsection{login している人のリスト: ¥texttt{who} コマンド}

 ¥texttt{who} は、現在、そのマシンにログインしている人のリストを表示する
コマンドだ。
 UNIX 系システムは、同時に何人もの人が使うことができるので、
このようなコマンドがある。
「¥texttt{who}」と打ってみよう。

macOSの場合:
¥begin{progls}
% ¥slbfl{who}
kokubo   console  Apr 25 09:15 
kokubo   ttys000  Apr 25 10:06
tomoda   ttys001  Apr 25 11:31   (192.168.1.5)
% 
¥end{progls}

WSL上のUbuntuの場合:
¥begin{progls}
$ ¥slbfl{who}
kokubo   pts/0        2025-04-25 09:15 (172.26.6.123)
tomoda   pts/1        2025-04-25 11:31 (172.26.6.124)
$ 
¥end{progls}

これが現在ログインしている人のリストである。

¥bigskip

それから、 ¥texttt{who} の仲間に、 ¥texttt{whoami} というコマンドがある。
これは、次のように「¥texttt{whoami}」と打つと現在使用中のユーザの ID
を表示する。

¥begin{progls}
% ¥slbfl{whoami}
kokubo
%
¥end{progls}


¥subsection{今日のニュースを取得: ¥texttt{curl} コマンド}

¥texttt{curl}は、URLからデータを取得するコマンドです。今日のニュースやヘッドラインを取得するために利用できます。
例えば、天気情報のAPIを利用して、東京の天気情報を取得してみましょう。

macOSの場合:
¥begin{progls}
% ¥slbfl{curl -s "https://wttr.in/Tokyo?format=%l:+%c+%t"}
Tokyo: 🌤 +23°C
%
¥end{progls}

WSL上のUbuntuの場合:
¥begin{progls}
$ ¥slbfl{curl -s "https://wttr.in/Tokyo?format=%l:+%c+%t"}
Tokyo: 🌤 +23°C
$
¥end{progls}

このようにインターネット上のリソースから情報を取得することができます。
「-s」オプションは、進捗情報など余分な出力を表示しないように指定するものです。

¥subsection{ランダムな名言を表示: ¥texttt{shuf} コマンドと組み合わせ}

現代のUNIX系システムでは、¥texttt{fortune}コマンドの代わりに、テキストファイルからランダムな行を選択する¥texttt{shuf}コマンドを使うことができます。あらかじめ名言を集めたファイルを作成しておけば、以下のようにして名言をランダムに表示することができます。

¥begin{progls}
% ¥slbfl{cat quotes.txt | shuf -n 1}
知識を得ることと知恵を得ることには違いがある。知識は事実やデータを集めること。知恵はそれをどう活かすかである。
%
¥end{progls}

あるいは、以下のようなシンプルな¥texttt{echo}コマンドを使って、ランダムな励ましメッセージを表示することもできます。

¥begin{progls}
% ¥slbfl{echo "今日も素晴らしい一日になりますように！"}
今日も素晴らしい一日になりますように！
%
¥end{progls}

これらのコマンドを組み合わせることで、古い¥texttt{fortune}コマンドと同様の機能を実現できます。

¥newpage

¥section{エラーからの回復方法}

人間は誰でも失敗することがある。
だから、コマンドを入力するときにも、いろいろ失敗するだろう。
それらの対処方法を改めて詳しく紹介しておこう。

¥begin{enumerate}
¥item ¥textbf{打ち間違えた文字を消したい}

打ち間違えた文字は、 ¥retrn を押す前なら、
 ¥KEYTOP{Delete} や ¥KEYTOP{Back Space} で消すことができる。

また、コマンドを途中まで打ったが、その行を全部キャンセルしたくなったら、
 ¥KEYTOP{Ctrl}+c と入力すればよい。

¥item ¥textbf{「Command not found.」(コマンドが見つからないよ。) と言われた}

これは、コマンドをタイプミスしたときにでるメッセージである。
よく見直してみよう。

¥item ¥textbf{コマンドを打ったら暴走した}

プログラムが暴走した場合には、 ¥KEYTOP{Ctrl}+c と打てば、たいてい止められる。

¥item ¥textbf{コマンドを打っても効かない}

たとえば X Window System でウィンドウを開くコマンドの後ろに「¥texttt{¥&}」を
付けないで実行すると、コマンドを実行したウィンドウには文字が入力できなくなる。
このときは、 ¥KEYTOP{Ctrl}+z と入力して、続けて「¥texttt{bg}」というコマンドを打てばよい。

¥item ¥textbf{コマンドを実行したが止め方がわからない}

ユーザの入力を対話的に処理するプログラムの多くは、
 ¥KEYTOP{Ctrl}+d で、終了できる。
また、ものによっては、「¥texttt{quit}」、「¥texttt{exit}」、「¥texttt{bye}」
などと入力すると、終了できることがある。

例を挙げよう。
¥texttt{bc} という電卓のプログラムがあって、
「¥texttt{bc}」と打つと起動できる。

¥begin{progls}
% ¥slbfl{bc}
bc 1.03 (Nov 2, 1994)
Copyright (C) 1991, 1992, 1993, 1994 Free Software Foundation, Inc.
This is free software with ABSOLUTELY NO WARRANTY.
For details type `warranty'. 
¥blsq
¥end{progls}
%¥newpage

後は、ここに数式を打つと、計算結果を表示してくれる。
例えば、「$1+2+3$」なら、

¥begin{progls}
% ¥slbfl{bc}
bc 1.03 (Nov 2, 1994)
Copyright (C) 1991, 1992, 1993, 1994 Free Software Foundation, Inc.
This is free software with ABSOLUTELY NO WARRANTY.
For details type `warranty'. 
¥slbfl{1+2+3}
6
¥blsq
¥end{progls}

このコマンドの実行を終了させるには、 ¥KEYTOP{Ctrl}+d、または「¥texttt{quit}」と入力する。

¥item ¥textbf{文字化けが起こった}

ターミナルアプリケーションを使っていると、
文字化けが起こることがある。

¥begin{figure}[H]
¥centering{¥includegraphics[scale=0.6]{mojibake.eps}}
¥end{figure}

このときは、以下の方法でリセットできます：

macOSのターミナルの場合：
「シェル」メニューから「ウィンドウをリセット」を選択するか、
 ¥KEYTOP{Command}+¥KEYTOP{K} キーでターミナルの内容をクリアします。

¥begin{figure}[H]
¥centering{¥includegraphics[scale=0.6]{terminal_reset_macos.eps}}
¥end{figure}

WSL上のUbuntuの場合：
「reset」コマンドを実行するか、¥KEYTOP{Ctrl}+¥KEYTOP{L} キーで画面をクリアします。

¥begin{progls}
$ ¥slbfl{reset}
$
¥end{progls}

また、SSHなどでリモート接続している場合には、
一度接続を切断して再接続するという方法もあります。
¥end{enumerate}

¥section*{この章で紹介したコマンド}

¥begin{center}
¥begin{minipage}{13cm}
¥begin{itembox}[l]{いろいろなコマンド}
¥begin{description}
¥item[¥texttt{uname} :] システム情報の表示

使い方: ¥texttt{uname} 「オプション」
¥item[¥texttt{date} :] 日時の表示

使い方: ¥texttt{date} 「フォーマット」
¥item[¥texttt{who} :] ユーザのリストの表示

使い方: ¥texttt{who}、¥texttt{whoami}
¥item[¥texttt{curl} :] ネットワークからデータを取得

使い方: ¥texttt{curl} 「URL」
¥item[¥texttt{shuf} :] ファイルからランダムな行を選択

使い方: ¥texttt{shuf -n 1} 「ファイル名」
¥end{description}
¥end{itembox}
¥end{minipage}
¥end{center}
%¥newpage
%{~}

